
\textit{Melius} ist eine innovative Plattform zur Erstellung und Verwaltung von digitalen Portfolios. Ziel des Projektes ist es, den Nutzer:innen eine einfache Möglichkeit zu bieten, ihre beruflichen Fähigkeiten, Erfahrungen und Projekte ansprechend und übersichtlich zu präsentieren.

Die Plattform besteht aus einem modernen, mit Angular entwickelten Frontend und einem leistungsfähigen, auf Java basierenden Backend mit dem Quarkus Framework. Für die Datenspeicherung wird eine PostgreSQL-Datenbank verwendet, die in einem Docker-Container läuft.

Die Benutzeroberfläche bietet eine zentrale Portfolioansicht, in der alle relevanten Informationen wie Lebenslauf, Kompetenzen und Projekte auf einer Seite zusammengefasst sind. Über Einstellungen können Sprache und Hintergrundfarbe angepasst werden, um das Portfolio individuell zu gestalten.

\subsection{Features für die Zukunft}

Als Erweiterung des bereits implementierten Funktionsumfangs wurden Überlegungen zu folgenden Funktionen angestellt:

\begin{itemize}

    \item 
    Die Möglichkeit, Standardvorlagen im Portfolio zu verwenden, die bereits ein Hintergrundbild, eine bestimmte Anordnung der Boxen oder ähnliches enthalten. Dadurch erhalten die Nutzer:innen eine Starthilfe und Inspiration beim Erstellen des Portfolios.

    \item 
    Die Möglichkeit, dem Portfolio konfigurierbare Formulare hinzuzufügen, bei denen die Überschrift und die Art der Eingabe festgelegt werden können. Dies soll zusätzliche Möglichkeiten bieten, die eigenen Stärken individuell darzustellen. Die verfügbaren Eingabearten entsprechen den bereits im Portfolio vorhandenen, wie z.B. Text, Slider oder Sternebewertung.

    \item 
    Für die Gruppenfunktion eine Erweiterung, die es Gruppenleiter:innen ermöglicht, in der Gruppenansicht Notizen an die Mitarbeitenden anzuhängen. Dies kann beispielsweise für Feedback oder Anmerkungen genutzt werden.

\end{itemize}
